% Generated draft: requires mathematical and source review.
\documentclass[11pt]{article}
\usepackage[margin=25mm]{geometry}
\usepackage{fontspec}
\setmainfont{Latin Modern Roman}
\usepackage{amsmath,amssymb,mathtools,unicode-math}
\setmathfont{Latin Modern Math}
\usepackage{enumitem,needspace,xurl,hyperref}
\setlength{\parindent}{0pt}
\setlength{\parskip}{5pt}
\setlength{\emergencystretch}{3em}
\newcommand{\N}{\mathbb{N}}
\newcommand{\Z}{\mathbb{Z}}
\newcommand{\Q}{\mathbb{Q}}
\newcommand{\R}{\mathbb{R}}
\newcommand{\C}{\mathbb{C}}
\newcommand{\F}{\mathbb{F}}
\newcommand{\A}{\mathbb{A}}
\newcommand{\PP}{\mathbb{P}}
\newcommand{\E}{\mathbb{E}}
\newcommand{\eps}{\varepsilon}
\newcommand{\dd}{\,\mathrm{d}}
\begin{document}
\section*{UnsolvedMath ID: 30001773}
\textbf{Source problem code:} catalogue entry
\subsection{Definitions and mathematical statement}
The supplied record collects several questions, rather than one fully specified conjecture. The following statements retain its separate targets and identify the missing qualifications.

\begin{enumerate}
\item Let \(F\) be a field, \(G\) a finite group, and \(H\leq G\). For each integer \(i\geq 1\), define
\[
T_i=\underbrace{FG\otimes_{FH}\cdots\otimes_{FH}FG}_{i\text{ factors}},
\]
with its \((FH,FH)\)-bimodule structure given by multiplication on the outer factors. Let \(\mathcal I_i\) be the set of isomorphism classes of indecomposable \((FH,FH)\)-bimodules occurring as direct summands of \(T_i\); multiplicities are discarded. Determine
\[
s_F(G,H)=\min\{i\geq 1:\mathcal I_i=\mathcal I_{i+1}\}.
\]
The minimum denotes a value only when the displayed set is nonempty. The record also requests the alternating-group case, but does not specify the subgroup inclusion intended in that case.

\item Write \(S_n\) for the symmetric group on \(n\) letters and \(A_n\) for its alternating subgroup. In characteristic \(p>0\), a vertex of a finite-dimensional indecomposable \(FK\)-module \(M\), for a finite group \(K\), is a \(p\)-subgroup \(P\leq K\), minimal under inclusion, for which \(M\) is a direct summand of \(\operatorname{Ind}_P^K U\) for some \(FP\)-module \(U\). A source at a vertex \(P\) is an indecomposable direct summand \(U\) of \(\operatorname{Res}_P^K M\) such that \(M\) is a direct summand of \(\operatorname{Ind}_P^K U\).

Determine the vertices and sources of the irreducible \(FS_n\)-modules denoted by \(D^\lambda\). In the usual split modular convention, \(\lambda\) is a \(p\)-regular partition of \(n\), meaning that no positive part occurs \(p\) or more times. The record does not explicitly specify the field hypotheses, characteristic range, or partition-label convention for this question; this convention therefore describes a conditional interpretation.

In characteristic \(2\), decide whether every nontrivial source of such an irreducible module has even \(F\)-dimension. If ``nontrivial'' means not isomorphic to the one-dimensional trivial \(FP\)-module \(F\), the proposed assertion is
\[
U\not\cong F\quad\Longrightarrow\quad \dim_F U\equiv 0\pmod 2
\]
for every source \(U\) at every vertex \(P\) of every module under consideration. The intended meaning of ``nontrivial'' is not explicitly defined in the record.

\item With the customary convention, \(\operatorname{Lie}_F(n)\) is the multilinear component of the free Lie algebra over \(F\) on \(x_1,\ldots,x_n\): it is spanned by Lie monomials in which every generator occurs exactly once. The group \(S_n\) acts by permuting the generators. For an \(FS_n\)-module \(M\), its dual \(M^*=\operatorname{Hom}_F(M,F)\) carries the action \((g\varphi)(m)=\varphi(g^{-1}m)\). A permutation module is an \(FS_n\)-module isomorphic to \(F[X]\) for a finite \(S_n\)-set \(X\).

The explicit Lie-module question in the supplied statement is whether
\[
\operatorname{Lie}_F(n)\otimes_F\operatorname{Lie}_F(n)^*
\]
is a permutation module under the diagonal \(S_n\)-action. The record does not specify restrictions on \(n\) or \(\operatorname{char}F\). Its original statement additionally requests Lie-module vertex, source, and indecomposability results, without identifying the exact modules, summands, hypotheses, or assertions for those additional questions. Those targets remain unresolved at the formulation level.

\item In characteristic \(2\), decide whether every nonsimple indecomposable Specht module for \(S_n\) remains indecomposable upon restriction to \(A_n\). Writing \(S_F^\lambda\) for the usual Specht module associated to a partition \(\lambda\vdash n\), the proposed assertion is
\[
S_F^\lambda\text{ indecomposable and not simple}
\quad\Longrightarrow\quad
\operatorname{Res}_{A_n}^{S_n}S_F^\lambda\text{ indecomposable}.
\]
The characteristic is specified, but no further hypotheses on \(F\) or restrictions on \(n\) and \(\lambda\) are supplied.

\item Determine the graded multiplicities of irreducible modules denoted by \(D^\lambda\) in the coinvariant algebras of symmetric groups and of \(\operatorname{GL}_n(q)\), or in their socles.

To describe the data needed for a precise version, for a specified finite group \(K\) acting linearly on a specified finite-dimensional \(F\)-space \(V\), one usual definition is
\[
C_F(K,V)=\operatorname{Sym}_F(V^*)\big/\left\langle\operatorname{Sym}_F(V^*)^K_{>0}\right\rangle,
\]
where the denominator is the ideal generated by homogeneous invariants of positive degree. For a finite-dimensional graded \(FK\)-module \(C=\bigoplus_{d\geq0}C_d\) and a simple module \(D\), its graded composition multiplicity is encoded by
\[
\sum_{d\geq0}[C_d:D]t^d.
\]
Here \([C_d:D]\) is the Jordan--H\"older multiplicity. The socle \(\operatorname{soc}_{FK}C\) is the sum of all simple submodules, and its graded multiplicities can be expressed using its degree components in the same way.

These definitions give a possible framework, not a determination of the intended coinvariant algebras. The record does not specify the representation \(V\), the coefficient field and its relation to the prime power \(q\), the labels for simple \(\operatorname{GL}_n(q)\)-modules, or whether ``multiplicities'' means composition multiplicities or another notion. It also leaves unclear whether both the whole algebra and its socle are requested, or whether these are alternative targets.
\end{enumerate}
\subsection{Short English statement}
This record bundles questions about tensor-chain stabilization, vertices and sources of symmetric-group modules, Lie modules, restriction of Specht modules to alternating groups, and graded multiplicities in coinvariant algebras. Several field assumptions and intended targets are missing. The supplied open status and verification claims are curation metadata; this statement-only draft requires human mathematical review.
\subsection{Sources}
\begin{itemize}
\item Supplied local UnsolvedMath registry record, permanent ID \(30001773\), title ``Vertices, Sources, and Filtrations of Symmetric-Group Modules,'' source code \texttt{catalogue entry}, in the Oberwolfach Reports Open Problems collection. Local registry references: \texttt{data/unsolvedmath/problems.json} and \texttt{catalog/problems.json}. The record gives proposed year \(2011\) and supplies no individual proposer.
\item Supplied report citation: \emph{Mini-Workshop: Modular Representations of Symmetric Groups and Related Objects} (2011), Oberwolfach Reports, pp. 1177--1205, DOI \texttt{10.4171/owr/2011/21}. Supplied source URL and literature-source link labelled ``original OWR report'': \url{https://doi.org/10.4171/owr/2011/21}. The supplied background describes Oberwolfach Reports as documenting mathematical workshops and their open-problem discussions. No problem-specific page or individual contribution is identified.
\item Supplied literature-source link labelled ``later Lie-module results'': \url{https://doi.org/10.1016/j.jalgebra.2015.04.015}. No authors, title, or theorem-level correspondence are supplied for this reference.
\item Literature assessment dated exactly \texttt{2026-08-21}: the supplied curation labels the bundled record open and states that later work gives special cases but not a full resolution. This is supplied dated curation, not independently verified. The background expressly requests independent verification of time-sensitive status and cited claims.
\item The supplied statement-status label is \texttt{corrected\_verified}; its verification note claims a check against the official Oberwolfach report and preservation of hypotheses, notation, and the open target. That claim is preserved as supplied metadata and is not an independent verification by this draft.
\end{itemize}
\end{document}
